\documentclass[10pt]{article}
\usepackage[margin=1in]{geometry}
\usepackage[T1]{fontenc}
\usepackage[utf8]{inputenc}
\usepackage{amsmath,amssymb}
\usepackage{graphicx}
\usepackage{booktabs}
\usepackage{array}
\usepackage[round,authoryear]{natbib}
\usepackage[hidelinks]{hyperref}
\urlstyle{same}
\emergencystretch=3em

\newcommand{\pending}[1]{\textbf{Results: pending} (produced by \nolinkurl{#1})}
\newcommand{\companion}[1]{\emph{#1}}
\newcolumntype{L}[1]{>{\raggedright\arraybackslash}p{#1}}

\title{Unequally Known: A Registered Report on a Tiered, Provenance-Tracked Global Synthetic Population for LLM Agents, Built from Synthetic Survey-Derived Demographics and Attitude Surveys}
\author{Real World Research\\\small Draft for internal review, 2026-10-11 (second revision).\\\small Stage-1 registered report: this paper reports no empirical results.}
\date{}

\begin{document}
\maketitle

\begin{abstract}
Simulations that use large language model (LLM) agents as stand-ins for people increasingly claim global scope, yet the data that ground those agents are far richer for some countries than for others, and the models studied so far answer cross-national opinion questions most like a small set of rich, mostly Western populations. This registered report specifies, before any data are touched, a resource that makes that inequality explicit: a synthetic population of about 100,000 agents over a country set fixed by a stated inclusion rule, built in two stages. A demographic skeleton is drawn from GLOPOP-S, a synthetic population built from household surveys, post-stratified to UN population totals; attitude vectors are attached by within-country donor matching from cross-national and regional attitude surveys. Every attribute of every agent carries its source, vintage and matching distance. We specify (i) an agent record schema with per-attribute provenance and a classification of items that refer to a specific state or period; (ii) four population tiers, written as an exhaustive decision list, and a per-country confidence score fixed before evaluation; (iii) noise-corrected diagnostics for cultural flattening and caricature; (iv) an uncertainty model in which a country's attitude estimates carry the donor survey's own sampling error, so that 100,000 agents are never reported as more informative than the surveys behind them; and (v) an evaluation protocol of eight hypotheses whose primary tests have uncertain outcomes, with a sample-size-matched noise ceiling under every error metric, targets fielded after the model's training cutoff where they exist, a simulated-borrowing arm, registered minimum sample sizes with a pre-registered downgrade to exploratory when they are not met, and a disclosure-risk gate on anything displayed or released. Agents are evaluated by their individual answers; the forecasting engine that consumes them queries the model per segment, and we state that gap rather than assume it away. A provisional reading of source terms (not legal advice) finds that every upstream source behind the open synthetic skeleton, and the largest attitude archive, restricts commercial use or redistribution; the build is therefore gated on written permissions, and agent-level attitude vectors will not be released. No construction or evaluation artifact exists yet; every result is marked pending with the artifact that will produce it.
\end{abstract}

\section{Introduction}

A live simulation of the world needs a population. The system this paper serves, Real World, keeps a population of synthetic agents, corrects its beliefs each morning against observed polls, with news as forcing, and publishes forecasts that are scored in public against prediction markets. Before any of that can mean anything, each agent needs to stand for someone: a person of a plausible age, sex, education, household and region in a real country, holding attitudes that people like them in that country actually hold.

Two facts make this hard. First, the ground truth is uneven. The largest harmonized archive of census microdata covers 104 countries and 656 censuses and surveys \citep{ipumshome}, but its terms bar commercial use \citep{ipumsterms}. The most complete open synthetic population is built from household surveys for 78 countries and fills 99 more, covering 15.7\% of world population, from the survey data of three countries chosen by human development index and geographic distance \citep{ton2024glopop}. Attitude surveys reach about 115 to 120 countries over four decades \citep{haerpfer2022wvs}; the Gallup World Poll, for example, typically interviews at least 1,000 people per country \citep{gallupwp}. Second, the models studied so far are biased in a known direction. In \citet{durmus2023}, default model answers to cross-national opinion questions were most similar to respondents from the USA, Canada, Australia and some European and South American countries; GPT-4o's answers sat closest to Finland, Andorra and the Netherlands \citep{tao2024}; alignment procedures created disparities across global opinions \citep{ryan2024}; and asking in the local language does not reliably fix it \citep{durmus2023,kwok2024,bulte2026}. These findings concern specific models; cultural prompting improved alignment for most countries in \citet{tao2024}, though it widened the gap for several European countries (for Finland, their distance measure rose from 0.20 to 2.43). A population presented as uniformly confident across countries nonetheless hides the countries where it is most likely to be wrong.

This paper specifies the population as a resource whose unevenness is a first-class output, and pre-registers the evaluation that will show whether that unevenness is measured usefully. It is written as a Stage-1 registered report \citep{nosek2018}: methods, hypotheses and analysis are fixed here; results come later from named artifacts.

\paragraph{Status.} A first engine exists (version 0.1.0, repository commit \texttt{ad95a99}), but the population this paper specifies does not: the engine samples $10^4$ agents for 16 countries from illustrative marginals written for the repository, not from GLOPOP-S or any survey (Section~\ref{sec:companions}). No number in this paper is a result of ours. Numbers about sources and prior work are quoted from, and cited to, the pages and papers we opened; a few numbers are design arithmetic (for example the floor's share of the sample, expected donors per cell, and power) and are labelled as such. Every experiment is a pre-registered protocol with a pending marker naming the artifact that will produce it (Table~\ref{tab:pending}).

\paragraph{Contributions.}
\begin{enumerate}
  \item \textbf{A provenance-first agent record.} Every demographic and attitude field records its source, wave, fieldwork year, donor country, match key, match distance and license row, and every item carries a pre-registered referent class (general, country-referential, time-referential) that governs whether a copied answer may be used at all (Sections~\ref{sec:schema} and~\ref{sec:items}).
  \item \textbf{Population tiers and a confidence score from inputs alone.} An ordered, exhaustive decision list assigns each country one of four tiers; a score in $(0,1]$ combines demographic source, donor recency, donor sample size and borrowing, with weights fixed here (Section~\ref{sec:tiers}). Model capability in the country's language is measured and recorded but kept out of the primary score, for reasons given in Section~\ref{sec:tiers}.
  \item \textbf{An honest uncertainty model.} Donor draws are proper multiple imputations (an approximate Bayesian bootstrap of the donor survey), so every released country estimate carries donor sampling error, not only agent sampling error (Section~\ref{sec:uncertainty}).
  \item \textbf{Noise-corrected flattening diagnostics.} A between-country slope of simulated on human deviations, corrected for measurement error, replaces an unstable ratio statistic; variance ratios subtract sampling noise; a default-ward shift is measured against the model's own unconditioned answers, net of sample-size bias (Section~\ref{sec:flattening}).
  \item \textbf{A pre-registered protocol} of eight hypotheses (Section~\ref{sec:protocol}) designed against the ways such evaluations pass trivially: primary tests that are magnitude tests or paired comparisons rather than tests the design guarantees, a noise ceiling matched to each comparison's sample sizes, post-cutoff targets against contamination with a registered feasibility gate, baselines that never see the target item, a simulated-borrowing arm, a test of the matching assumption on variables (including region) deliberately left out of the key, and a non-inferiority test of the released object against the donor surveys themselves.
\end{enumerate}
Section~\ref{sec:sources} also reports a provisional reading of the terms of fifteen sources. It is not a contribution and not legal advice; it explains why the build is gated and the release is thin.

\paragraph{What is new, and what is not.} Population synthesis \citep{gallagher2018spew}, statistical matching \citep{dorazio2006,rassler2002}, datasheets \citep{gebru2021} and value-profile personas \citep{abels2026} each exist. SPEW builds world populations without attitudes; \citet{abels2026} sample value profiles from target-country survey responses but do not attach them to a census-calibrated skeleton, track provenance per field, or say how well each country is known. The contribution here is the combination and its evaluation: a global skeleton joined to attitude donors with provenance on every field, a per-country grade fixed before any outcome is seen, and a registered test of whether that grade predicts error.

\paragraph{What this paper does not claim.} It does not claim the agents are accurate for any country; that is what the protocol tests. It does not claim that grounding removes Western bias; the literature suggests it reduces but does not remove it \citep{tao2024,bulte2026}. It does not claim individual-level fidelity: an agent is a draw from a conditional distribution, not a person. It does not claim that 100,000 agents carry more attitude information than the few thousand survey respondents behind them; they carry less. And it does not validate the agents on what the product mostly asks of them, namely novel event questions and daily belief updates, nor the segment-level path the forecasting engine uses; Section~\ref{sec:limits} states what the hypotheses cover.

\section{Related work}
\label{sec:related}

\paragraph{Population synthesis.} Iterative proportional fitting rescales a seed table until it matches known margins \citep{deming1940}. Integer populations can be obtained from fractional weights by truncate--replicate--sample \citep{lovelace2013trs}. In a benchmark of four synthetic reconstruction methods on French municipalities, hierarchical IPF and entropy minimisation performed best and the integerisation step mattered \citep{yameogo2021}. Deep generative models can produce agents not present in the seed sample \citep{borysov2019}. SPEW is an open-source framework for generating synthetic populations and ecosystems for countries worldwide \citep{gallagher2018spew}. None of these attaches attitudes; we use the simplest skeleton method and spend the effort on the attitude layer.

\paragraph{Global microdata and synthetic populations.} IPUMS International harmonizes census samples contributed by national statistical offices \citep{ipums2025,ipumshome}. GLOPOP-S provides about 7.3 billion synthetic individuals for 2015 with age, education, gender, income or wealth, settlement type and household structure, built from LIS and DHS household surveys; for each of 99 countries without survey data (15.7\% of world population) it takes the ten countries with the most similar Global Data Lab human development index and combines the survey data of the three geographically closest, and about 1\% of world population has no synthetic file \citep{ton2024glopop}. Global Data Lab publishes subnational development indicators \citep{smits2019shdi}. UN World Population Prospects 2024 (WPP) gives estimates and projections for 237 countries or areas \citep{wpp2024}. None of these carries attitudes.

\paragraph{Attitude surveys.} The joint WVS/EVS trend page states 452 surveys from 115 countries or territories in its text and lists 120 countries and 666,907 cases in its table, over 1981--2022 \citep{haerpfer2022wvs,evs2022trend}; we report both because the source does. Regional barometers add depth: Afrobarometer lists 45 countries \citep{afrobarometer}; Arab Barometer lists nine waves since 2006, of which eight have released data sets \citep{arabbarometer}; Latinobarómetro runs about 20,000 interviews a year in 18 countries \citep{latinobarometro}; and the Asian Barometer is releasing its sixth wave country by country \citep{asianbarometer}. The Gallup World Poll typically interviews at least 1,000 people aged 15 and over per country across more than 160 countries \citep{gallupwp}, and Pew's Global Attitudes surveys covered between 19 and 35 countries per spring wave from 2022 to 2025 \citep{pewdatasets}. The cross-cultural map of \citet{inglehart2005} organizes WVS societies along two value dimensions.

\paragraph{LLMs as survey respondents.} Conditioning a model on demographic backstories can reproduce US subgroup opinion patterns \citep{argyle2023,moon2024}, but synthetic samples misstate variance and regression relationships and are unstable over time \citep{bisbee2024}, model opinions are unrepresentative of many US groups \citep{santurkar2023}, and model survey answers are driven by option ordering and labelling \citep{dominguez2024} and show response biases unlike humans' \citep{tjuatja2024}. In a WVS microdata study of US respondents, multi-attribute persona prompting gave no clear aggregate improvement and degraded some subgroups \citep{tadaymorocho2026}. Agents grounded in two-hour interviews with 1,052 Americans reached 83\% of participants' own two-week test--retest consistency on General Social Survey items; agents given each participant's own survey answers reached 82\%, the two combined 86\%, and demographics-only agents 74\% \citep{park2024}; Twin-2K-500 releases 500-question answers from over 2,000 people for such individual-level work \citep{toubia2025twin}. Those agents received the same person's answers. Donor matching gives an agent another respondent's answers, matched on four demographics, so Park et al.'s 82\% does not transfer to our design; whether donor vectors carry useful signal is what H3 tests. Pew's 2026 study conditioned Claude Opus 4.6 on each American Trends Panel member's demographics, 71 political typology variables and an expert-reflection step; synthetic estimates missed human ones by 11.3, 14.7 and 11.2 percentage points on three replicated waves, 12.4 on average (the methodology page gives 11.4 for the development wave), and in a comparison on a subset of questions GPT-5.1 described a more extreme public and Opus a more middle-of-the-road one than the humans \citep{pew2026silicon}. A Korean synthetic panel showed 14--19 point mean absolute error while ranking items reasonably (correlations 0.70--0.91), and calibrating on real data did not beat direct estimation from the real subsample \citep{kim2026korean}. \citet{ku2026} give a model one respondent's answers on one item set of a Taiwanese election survey and predict that same respondent's answers on a disjoint set; H3 has the same held-out-item structure but conditions on a matched donor's answers. \citet{abels2026} build personas from value profiles sampled from target-country survey responses and calibrate response diversity, the closest prior to our Stage~2.

\paragraph{Describing versus sampling distributions.} Instruction-tuned models can describe a response distribution in one call more accurately than they can sample it across personas \citep{jang2026,meister2025}; this makes our direct-distribution baseline (B3) strong and predicts within-country variance collapse. Independent persona agents grounded on Turkish WVS respondents piled onto a modal answer \citep{ozkan2026}. Matching published marginals does not establish respondent simulation: a no-persona population query beat committed synthetic panels on aligned batteries \citep{doudkin2026}, which is why H8 compares the released population with the donor surveys themselves rather than only with toplines. Fine-tuning on survey distributions is a further strong non-persona approach \citep{suh2025subpop}; we do not fine-tune. Replacing participants with LLMs can misportray and flatten identity groups \citep{wang2024flatten}. Benchmarks include SimBench \citep{hu2026simbench}, WorldValuesBench \citep{zhao2024wvb}, GlobalOpinionQA \citep{durmus2023} and BLEnD for everyday cultural knowledge \citep{myung2024blend}; the WVS-derived ones are built from public items, which is why Section~\ref{sec:protocol} treats contamination as a design problem. Stated training cutoffs need not match a model's effective cutoff for a given source \citep{cheng2024dated}.

\paragraph{Large agent populations.} AgentSociety simulates over 10,000 LLM-driven agents \citep{piao2025agentsociety}, and LLM archetypes let one model call stand for many agents in a simulation of 8.4 million agents \citep{chopra2024limits}, the per-segment pattern our forecasting engine also uses (Section~\ref{sec:companions}). Persona Hub curates one billion personas from web text \citep{ge2024personas}; those personas are not calibrated to any population, which is the alternative this resource is built against.

\paragraph{Cultural bias and language.} \citet{henrich2010} named the WEIRD (Western, educated, industrialized, rich, democratic) bias of behavioural science. \citet{atari2023} report two findings: LLM responses to psychological measures are an outlier relative to large-scale cross-cultural data, and LLM performance on cognitive psychological tasks most resembles that of people from WEIRD societies and declines with cultural distance from them ($r=-.70$). Evidence on prompt language conflicts: translated prompts did not move answers toward the speakers of that language in \citet{durmus2023}; native-language prompting reduced alignment in a 15-country persona simulation \citep{kwok2024}; Arabic prompting improved alignment with Egypt \citep{alkhamissi2024}; matching prompt language to the questionnaire language steered models toward the corresponding country in \citet{liu2025}; explicit cultural framing beat prompt language, and combining them added nothing, across 10 models and 11 languages \citep{bulte2026}; multilingual ability does not imply cultural alignment \citep{rystrom2025}; and non-English questions do not reliably remove English-culture answers \citep{wang2023thanksgiving}. Non-English text also costs more: tokenizers can produce up to 15 times as many tokens for the same content across languages \citep{petrov2023}, and commercial API cost varies by language with tokenization \citep{ahia2023}.

\paragraph{Small-area estimation and imputation.} Multilevel regression and poststratification estimates subgroup opinion from a model fitted on respondents and reweighted to population cells \citep{park2004mrp}, including deep interactions among demographic variables \citep{ghitza2013}; we use it as a non-LLM baseline (B6). Hot-deck imputation and its variance are reviewed by \citet{andridge2010}; proper multiple imputation follows \citet{rubin1987}.

\paragraph{Dataset documentation.} We follow datasheets for datasets \citep{gebru2021}, extended with per-attribute provenance and per-country confidence (Appendix~\ref{app:datasheet}).

\section{Sources and a provisional terms reading}
\label{sec:sources}

Table~\ref{tab:sources} lists each source, its planned role, and its terms as we read them on 2026-10-11. This is a reading by non-lawyers, not legal advice. Three findings shape the design.

\paragraph{The open skeleton rests on restricted upstreams.} GLOPOP-S is released under CC BY 4.0 \citep{glopopdataverse}, but it is built from LIS, whose microdata access is ``exclusively for non-commercial purposes'' \citep{lisaccess}, from DHS, whose terms require that data ``will not be used for any marketing or commercial venture'' and that micro-level data ``will not be re-distributed'' \citep{dhsterms}, and it selects fill countries with Global Data Lab indicators, whose commercial use ``is only possible after written permission'' \citep{gdlterms}; the Global Data Lab Area Database itself is no longer offered and has been split into thematic datasets \citep{gdlareadata}. Whether the CC BY label settles commercial reuse of the derived synthetic file is a question we have put to counsel and to the GLOPOP-S authors. IPUMS International states that commercial use ``is strictly prohibited'' and bars redistribution \citep{ipumsterms}; we do not sample from it, and use it only for aggregate-only validation if IPUMS grants written permission (H1).

\paragraph{The main attitude archive is non-profit only.} The WVS trend-file download requires agreeing to conditions under which files are available ``provided: a) that they are used for non-profit purposes'' and ``the data files themselves are not redistributed'' \citep{wvsconditions}, and the file's registered metadata state ``No commercial usage'' \citep{haerpfer2022wvs}. Real World is a commercial project, so WVS/EVS donors cannot be used, even for the evaluation run, without written permission. Afrobarometer states that its data ``are protected by copyright'' and requires acknowledgement, and is silent on commercial use \citep{afropolicy}; a copyright notice without a grant is not a permission. Arab Barometer's linked policy covers reprints of its publications and says the data download is free, and is silent on commercial use of the data \citep{arabpolicy}. We found no terms for Latinobarómetro. Pew licenses datasets non-exclusively, forbids publication in full or substantially in full, and forbids attempts to identify respondents \citep{pewterms}; we use Pew only as published aggregates, never as donors. GlobalOpinionQA is CC BY-NC-SA 4.0 \citep{globalopinionshf}. The multilingual item set used to measure $q^{\text{lang}}$ is MIT-licensed \citep{mmmlu}.

\paragraph{Consequence: the build fails closed.} Each source has a row in \nolinkurl{licenses.csv} holding the written answer, if any. The loader refuses any source whose row does not read \texttt{permitted} for the intended use (construction, evaluation, or release). Until written permissions arrive, no part of this protocol runs on WVS/EVS, DHS-derived, LIS-derived or GDL-derived data in a commercial setting. An academic collaborator holding their own licences is not a way around this: a protocol run for the benefit of a commercial company with a declared financial interest may fall outside non-commercial or non-profit terms even when an academic runs it. That route is used only for a source whose custodian confirms in writing that this use, with this beneficiary, is within its licence; the confirmation goes in \nolinkurl{licenses.csv}, and we will say which route was used for each source.

\begin{table}[!htbp]
\centering
\small
\caption{Sources as stated by each source on 2026-10-11. ``Role'' is our planned use. ``Release'' says what derived from the source can appear in the public resource; no source permits agent-level attitude release under this design (Section~\ref{sec:release}).}
\label{tab:sources}
\begin{tabular}{@{}L{2.75cm}L{3.1cm}L{2.0cm}L{4.3cm}L{1.6cm}@{}}
\toprule
Source & Coverage (as stated) & Role & Terms as read & Release \\
\midrule
IPUMS International & 104 countries, 656 censuses/surveys & Validation only, if permitted & Commercial use prohibited; no redistribution & No \\
GLOPOP-S & 78 countries from LIS/DHS; 99 filled from 3 similar countries; base 2015 & Demographic skeleton & CC BY 4.0; upstream terms below & Pending counsel \\
LIS (upstream) & -- & Via GLOPOP-S & Non-commercial access only & No \\
DHS (upstream) & 45 of GLOPOP-S's 78 survey countries & Via GLOPOP-S & No commercial venture; no redistribution; no sharing without written consent & No \\
Global Data Lab & Subnational indicators & Via GLOPOP-S fill & Non-commercial; commercial use needs written permission & No \\
UN WPP 2024 & 237 countries or areas & Post-stratification totals; migration for H1 & CC BY 3.0 IGO shown on a WPP 2024 brief \citep{wpp2024brief}; data download page not confirmed & Margins only \\
WVS/EVS trend & 115 (text) or 120 (table) countries/territories; 666,907 cases; 1981--2022 & Attitude donors & Non-profit use; files not redistributed; metadata ``No commercial usage'' & No \\
Afrobarometer & 45 countries & Donors, targets & Copyright notice; attribution required; silent on commercial use & No \\
Arab Barometer & 9 waves listed, 8 with data sets, 2006--2026 & Donors, targets & Download free; linked policy covers reprints only & No \\
Latinobarómetro & 18 countries, $\approx$20,000 interviews/yr & Donors, targets & No terms found & No \\
Asian Barometer & Wave 6 releases ongoing & Not used & Access terms not described & No \\
Gallup World Poll & 160+ countries, usually $\geq$1,000 per country, ages 15+ & Aggregate targets (H8) & Subscription or purchase & Aggregates if licensed \\
Pew Global Attitudes & 19--35 countries per wave & Aggregate targets (H8) & Excerpts only; no re-identification & Aggregates only \\
GlobalOpinionQA & 2,556 questions & Not used (contamination reference) & CC BY-NC-SA 4.0 & No \\
MMMLU & MMLU test set in 14 languages, human-translated & $q^{\text{lang}}$ items & MIT & Item ids only \\
\bottomrule
\end{tabular}
\end{table}

\section{Resource design}
\label{sec:design}

\subsection{Country set, frame and allocation}
\label{sec:alloc}

\paragraph{Notation.} $n_k$ is the number of agents for country $k$; $m_k$ is the number of completed interviews in $k$'s donor wave; $n^{\text{tgt}}$ is the size of an evaluation target sample; $n^{\text{eval}}$ the number of evaluation agents. These are kept distinct because H4 controls for target size while the score uses donor size.

\paragraph{Inclusion rule.} The country set contains every WPP 2024 country or area \citep{wpp2024} whose estimated population in the agent frame in the target year $t = 2026$ is at least 100,000. $C$, the number of included countries, is computed from WPP before any survey file is opened and committed in \nolinkurl{coverage.csv}; it is at most 237. The country count is therefore an output of a fixed rule, not a target.

\paragraph{Frame.} A country's agent frame is the age frame of the survey that supplies its attitudes: its own donor survey, or for a borrowed country the fallback country's donor survey (Section~\ref{sec:fallback}); 18 and over for most programs, 15 and over where that survey's frame starts at 15. Tier~4 countries (no donor survey) use 18 and over. Agents are therefore never younger than the survey that supplies their attitudes. GLOPOP-S ages come in bands (0--4, 5--14, 15--24, then ten-year bands to 65+) \citep{ton2024glopop}. One single-year rule is used everywhere: each record receives a single year of age drawn from the WPP single-year age--sex distribution within its band for the skeleton's reference year (2015 before any cohort shift, Section~\ref{sec:stage1}); records below the frame are then dropped.

\paragraph{Allocation.} With population $P_k$ in the frame, a total of $N$ agents and a floor $f$ per country,
\begin{equation}
n_k = f + \operatorname{LR}\!\left[(N - fC)\,\frac{P_k}{\sum_j P_j}\right],\qquad w_a = \frac{P_{k(a)}}{n_{k(a)}},
\end{equation}
where LR is largest-remainder rounding so that $\sum_k n_k = N$. We pre-register $N = 100{,}000$ and $f = 100$. Since $C \le 237$, the floor uses at most 23,700 agents (design arithmetic). Every world or regional aggregate must be computed with $w_a$, and the release ships the weights. Tier~4 countries carry no attitude block (Section~\ref{sec:tiers}), so every world or regional attitude aggregate is computed over Tiers 1--3 with weights renormalized over the included countries, describes the included population only, and is published with the share of the frame population it excludes.

A floor of 100 agents bounds only the agent-sampling part of a country proportion's standard error, at $\sqrt{0.25/100} = 0.05$ (design arithmetic). The donor part does not shrink as $n_k$ grows: for a donor wave of $m_k \approx 1{,}000$ respondents with design effect $d$, it is about $\sqrt{0.25\,d/1000} \approx 0.016\sqrt{d}$ however many agents are drawn (design arithmetic). A large country's $n_k$ of order $10^4$ agents draws on about $10^3$ donors, so each donor vector is reused ten or more times. The attitude block of 100,000 agents carries no more information than the donor files behind it, and the release says so; Section~\ref{sec:uncertainty} gives the standard errors it ships.

\subsection{Agent record and provenance}
\label{sec:schema}

An agent record has three blocks. The \emph{demographic block} holds country (ISO 3166-1 alpha-3), subnational region as defined by GLOPOP-S (Global Data Lab subnational regions, which are not always first-level administrative units), settlement type (urban or rural), single-year age, sex, education level, income or wealth quintile (LIS income or DHS wealth index, which are not comparable across countries \citep{ton2024glopop}), household size and type, and relationship to household head. The \emph{attitude block} holds a vector of harmonized items with missing-value codes and referent flags (Section~\ref{sec:items}). The \emph{runtime block} holds the prompt language, the model and version, and an assimilation state managed elsewhere.

Every field $f$ of agent $a$ carries a provenance tuple
\begin{equation}
\pi_{a,f} = (\text{source},\ \text{wave},\ \text{fieldwork year},\ \text{donor country},\ \text{key},\ \text{distance},\ \text{collapses},\ \text{license}),
\end{equation}
where key and distance are the match key and match distance, collapses lists the imputation-class collapse steps, and license points to the source's row in \nolinkurl{licenses.csv}. Every agent carries its design weight $w_a$ and its country's tier and confidence components. Provenance is written when the field is filled, not inferred later.

\subsection{Stage 1: demographic skeleton}
\label{sec:stage1}

For each country $k$ with a GLOPOP-S file, the seed is its synthetic individuals in the agent frame. GLOPOP-S describes 2015, so we post-stratify the seed to the WPP age--sex totals for $t$ \citep{wpp2024}. With one joint age--sex margin, raking reduces to post-stratification; the joint distribution of the remaining attributes given age and sex is kept from 2015.

Plain reweighting keeps 2015's relation between age and education, which is biased under educational expansion: the people aged 35--44 in 2026 were mostly aged 25--34 in 2015 and kept that cohort's education. We therefore pre-register a \emph{cohort-shift} alternative for target years after 2015. Each 2015 record, with its single year of age from the rule in Section~\ref{sec:alloc}, is aged by $t - 2015$ years, so cohorts keep their education. Two corrections apply to the young. Records aged 15--24 in 2015 were often still in education; after shifting, any such record whose new age is 25 or more has its education redrawn from the 2015 distribution of the 25--34 band with the same sex and settlement type, and never lowered below its 2015 level. This assumes completion by 25 and that the cohort reaches at least the education of the cohort ten years older, which understates attainment where expansion continued; we state the direction. New entrants to the frame take the 2015 education distribution of the band they enter at the same age. Differential mortality by education is ignored. The result is post-stratified to WPP age--sex totals. Cohort shift is defined only forward in time; for a census year before 2015 only plain reweighting is evaluated. Plain reweighting is the primary build; the H1 primary test decides, by the rule in Table~\ref{tab:tests}, whether a new, versioned release switches to cohort shift.

For the countries without a GLOPOP-S file (about 1\% of world population \citep{ton2024glopop}), the skeleton is WPP-only: age and sex from WPP, all other demographic fields coded missing, and $q^{\text{demo}}_k = 0.25$.

We draw $n_k$ agents by systematic probability-proportional-to-weight sampling, stratified by region and settlement type. We draw rather than materialize the full integer population, which would need an integerisation step such as truncate--replicate--sample \citep{lovelace2013trs}. \nolinkurl{coverage.csv} reports, per country, the distribution of the key $X$ in the skeleton and in the donor survey, since the agents' attitude aggregate is in effect the donor survey reweighted to the skeleton's $X$.

\subsection{Stage 2: attaching attitudes by donor matching}
\label{sec:stage2}

\paragraph{Eligible wave.} One definition is used everywhere (Stage~2, tiers, fallback, H7). A wave of any attitude program is \emph{eligible} for country $k$ if (i) its fieldwork midpoint $t_k$ satisfies $t - t_k \le 15$ years; (ii) it has at least 500 completed interviews in $k$ from a national sampling frame; (iii) it records sex, age, settlement type and education in a form the committed crosswalk maps; (iv) it contains at least half of the common item set $J_c$ (Section~\ref{sec:items}); and (v) its source's license row reads \texttt{permitted}. The \emph{donor wave} for $k$ is the eligible own-country wave with the latest fieldwork midpoint, ties broken by larger $m_k$. A wave is \emph{recent} if $t - t_k \le 8$ and $m_k \ge 1{,}000$.

\paragraph{Matching.} The attitude block is filled by statistical matching \citep{dorazio2006,rassler2002}: file A (the skeleton) has demographics $X$; file B (the donor wave) has $X$ and attitudes $Y$. The key is $X = $ (sex, age band, settlement type, education). Education is crosswalked from the five GLOPOP-S levels (less than primary, complete primary, incomplete secondary, ``complete secondary or tertiary'', higher \citep{ton2024glopop}) to ISCED 2011 levels \citep{isced2011} by a committed table. Because GLOPOP-S level 4 spans secondary and some tertiary, a skeleton level maps to a \emph{set} of ISCED levels, and the education distance between an agent and a donor is zero if the donor's level lies in the set and otherwise the ordinal distance to the nearest member. Crosswalk ambiguity is therefore part of the match distance. Income or wealth is not in the key, because its skeleton coding differs by country. Region is not in the key: with the cell sizes below, adding it would collapse most cells; its omission is tested in H2a and stated in Section~\ref{sec:limits}.

\paragraph{Expected cell sizes (design arithmetic).} The key has up to $2 \times 6 \times 2 \times 5 = 120$ cells (sex, six frame age bands, settlement type, five education levels). A donor wave of about 1,000 respondents leaves about 800 donors after the 80/20 evaluation split, about 7 per cell on average, and far fewer in rare cells. Collapse is therefore common, not exceptional; \nolinkurl{coverage.csv} reports the share of agents matched at each collapse step.

\paragraph{Imputation classes and draws.} For agent $a$ in country $k$, the imputation class is the respondents of the donor file matching exactly on sex, age band and settlement type; among them, a donor is drawn with probability proportional to survey weight among those at minimum education distance. If the class has fewer than ten respondents, age bands are collapsed, then settlement type, each step recorded in $\pi$. The agent receives the donor's item vector, which keeps within-person correlations among items, subject to the referent rules of Section~\ref{sec:items}. We store $M$ donor draws per agent. Each draw is a proper imputation in Rubin's sense \citep{rubin1987,andridge2010}: draw $m$ first resamples the donor wave (primary sampling units with replacement within strata where the file identifies them, the program's $m$-th replicate weight set where it supplies replicate weights, respondents otherwise), then hot-decks from that resampled file. $M = 5$ for model-answered evaluation items and $M = 20$ for released aggregates of donor-carried items, which need no model calls. Draw $m=1$ is the default agent state.

\paragraph{The matching assumption.} Matching rests on conditional independence: given $X$, attitudes are independent of any skeleton attribute not in $X$ (region, household type, income or wealth). It cannot be tested between files \citep{dorazio2006}, but it can be probed inside one file: H2a hides survey variables, including region, from the key and measures how much of their association with attitudes the match recovers. We also publish, for each pair of a dichotomized attitude event $y$ and a dichotomized skeleton-only event $z$ used downstream (including region indicators), the Fréchet bounds
\begin{equation}
\sum_x p(x)\max\{0,\ p(y\mid x)+p(z\mid x)-1\} \le p(y,z) \le \sum_x p(x)\min\{p(y\mid x),\ p(z\mid x)\},
\end{equation}
with bootstrap intervals on the widths \citep{dorazio2006,statmatch}. The width says how much any cross-tabulation of $y$ by $z$ is an artifact of the assumption. Widths are descriptive: a holdout joint distribution from the same population lies inside its own Fréchet bounds by construction, so coverage of the bounds is not used as a test.

\subsection{Fallback donors}
\label{sec:fallback}

When $k$ has no eligible own-country wave, donors come from one fallback country $k'$ chosen by a rule that uses only the skeleton and WPP, so it needs no further licensed source. Candidates are the countries with an eligible own wave in the same UN M49 sub-region \citep{m49}; if there are none, the same M49 region. Among candidates, $k'$ minimizes the Euclidean distance between standardized country vectors computed from the post-stratified skeletons: median age, urban share, and the shares in GLOPOP-S education levels 1--2 and 4--5. For WPP-only countries the vector is median age alone. Ties go to the candidate with the larger donor wave, then to alphabetical ISO3 order. We do not reuse GLOPOP-S's own fill countries, because each filled country was built from three countries, not one \citep{ton2024glopop}. The map is computed from metadata before any outcome is seen, committed as \nolinkurl{fallback_map.csv}, and every borrowed field is marked. Demographic distance is a weak proxy for attitudinal similarity (within one sub-region it can pair very different societies), so the simulated-borrowing arm also evaluates one alternative rule as a secondary comparison: the nearest candidate in the same M49 region by position on the two dimensions of \citet{inglehart2005}, computed from published country aggregates. Borrowing asserts that a country thinks like its neighbour; it is penalized in the score and measured directly by the simulated-borrowing arm (Section~\ref{sec:protocol}).

\subsection{Items, scales and referents}
\label{sec:items}

\paragraph{Common item set.} Within one program, every item fielded in a wave is common across that wave's countries with no crosswalk. Across programs, the \emph{common item set} $J_c$ contains items that satisfy three rules, applied to questionnaires only: (i) a committed crosswalk maps the item across programs with the same stem construct and a one-to-one mapping of response categories; (ii) the item is fielded in eligible waves of at least 20 countries; (iii) its referent class (below) is general. Items without a one-to-one mapping stay in country-specific vectors but are excluded from every cross-program statistic. All cross-country statistics (H4, H5, H8) use $J_c$ or a single program's common items, so they compare the same items on the same scale. Distributional distances are computed only between distributions on the same category set. ``Don't know'' and refusal are kept as one category, the same way for humans and agents, wherever the questionnaire offers it. Means are taken on ordinal items scored at equally spaced positions in $[0,1]$, excluding the don't-know category; nominal items enter mean-based statistics only as per-category indicator shares.

\paragraph{Referent classes.} An answer copied from another person is valid only if the question means the same thing to the agent. Each item is classified before any run as \emph{country-referential} (it names or refers to the respondent's own state, government, leaders, national institutions, national economy or the country's most important problem), \emph{time-referential} (it refers to a current or recent period, event or incumbent), or \emph{general}. Two coders classify independently from the questionnaire text; disagreements go to the stricter class; the classification and inter-coder agreement are committed in the registry. In a \emph{borrowed} vector, country-referential items are dropped (coded not applicable); time-referential items are dropped too. In an own-country vector from a wave that is not recent, time-referential items are dropped and country-referential items are flagged \emph{stale-referent} and excluded from released aggregates and from product displays. A dropped item has no donor answer; it is answered by the model only where a hypothesis says so, and such answers are flagged as model-generated.

\subsection{Language assignment}
\label{sec:lang}

Each agent has a prompt language. The default is English with an explicit cultural frame (country, region, the agent's own attributes and donor answers rendered as a profile), because explicit framing outperformed prompt language in the broadest comparison we found \citep{bulte2026} and native-language prompting lowered alignment in a 15-country simulation \citep{kwok2024}. The evidence is mixed \citep{alkhamissi2024,liu2025} and none of it tests our model, so H6 tests the alternative in 20 countries committed before the run. Native language can only ever be enabled for those 20 countries, and only if H6 passes for that country; it is not part of the launch configuration.

Model capability is uneven across languages: a vendor table reports Claude Haiku 4.5 with extended thinking at 96.4\% of its English score in Spanish but 52.7\% in Yoruba on MMLU translated into 14 languages by professional translators \citep{anthropicmultilingual}. That is not the model our run uses, so we measure $q^{\text{lang}}$ ourselves (Section~\ref{sec:tiers}). Token cost differs by language \citep{petrov2023,ahia2023}, so H6 records measured token counts on the run model's own tokenizer.

\subsection{Population tiers and per-country confidence}
\label{sec:tiers}

\paragraph{Tiers.} Let $D_k \in \{\text{own}, \text{fill}, \text{wpp}\}$ be the demographic source (an own-country GLOPOP-S survey file, a GLOPOP-S fill, or WPP only). Let $A_k$ be \emph{own-recent} if $k$'s donor wave is recent, \emph{own-other} if $k$ has a donor wave that is not recent, \emph{borrowed} if $k$ has no eligible own wave but its fallback $k'$ has one, and \emph{none} otherwise. The tier is the first rule that applies:
\begin{enumerate}
  \item $A_k = \text{none}$ $\Rightarrow$ \textbf{Tier 4, ungrounded.}
  \item $A_k = \text{borrowed}$ or $D_k \in \{\text{fill}, \text{wpp}\}$ $\Rightarrow$ \textbf{Tier 3, borrowed} (attitudes, demographics, or both).
  \item $A_k = \text{own-recent}$ (and $D_k = \text{own}$) $\Rightarrow$ \textbf{Tier 1, grounded.}
  \item otherwise ($A_k = \text{own-other}$, $D_k = \text{own}$) $\Rightarrow$ \textbf{Tier 2, partly grounded.}
\end{enumerate}
The list is exhaustive and its outcomes are mutually exclusive. An own-country wave older than 15 years is not eligible, so such a country is Tier~3 if its fallback has an eligible wave and Tier~4 otherwise. Tier~3 mixes two cases: a country with borrowed attitudes, and a country with its own recent attitudes but filled demographics. The second can have a higher $c_k$ than a Tier~2 country; tiers and scores are separate signals and both are shown, and the tier records which component is borrowed. Tier~4 agents carry the demographic block only: no attitude block is built and the model is not used to invent one. They are shown and flagged, have no confidence score, and are excluded from attitude aggregates (Section~\ref{sec:alloc}).

\paragraph{Confidence score.} For a country in Tiers 1--3,
\begin{equation}
c_k = (q^{\text{demo}}_k)^{1/3}\,(q^{\text{att}}_k)^{2/3},
\end{equation}
a weighted geometric mean, so a weak component pulls the score down. The components are:
$q^{\text{demo}}_k = 1$ for $D_k=\text{own}$, $0.5$ for a fill, $0.25$ for WPP only;
$q^{\text{att}}_k = b_k\, e^{-(t-t_k)/\tau}\min(1, m_k/1000)$ with $t_k$ the fieldwork midpoint of the donor wave used, $m_k$ its size, $\tau = 8$ years, and $b_k = 0.5$ for borrowed donors, $1$ otherwise. Every component is strictly positive for a scored country. In deployment $t = 2026$; in H4, $c_k$ is recomputed with $t$ set to the fieldwork midpoint of the target the error is measured against, so the recency factor refers to the same gap as the error it is meant to predict.

\paragraph{Why language capability is out of the primary score.} An earlier draft multiplied in $q^{\text{lang}}_k$, the run model's accuracy in country $k$'s language relative to English. We keep it out of $c_k$ for three reasons. The launch configuration prompts in English, so $q^{\text{lang}}$ measures capability in a language the system does not use; its only rationale is as a proxy for how much the model knows about the culture, and multiple-choice knowledge accuracy is a weak proxy for opinion fidelity. The language rule (the language most respondents of $k$'s latest wave were interviewed in) assigns English or French to much of anglophone and francophone Africa and South Asia, so $q^{\text{lang}} \approx 1$ exactly where cultural knowledge is likely weakest. And the human-translated item set covers 14 non-English languages \citep{mmmlu}; with a worst-case value for every other language, the component becomes an indicator of membership in that list, collinear with development level. $q^{\text{lang}}_k$ is still measured (200 committed MMMLU items, accuracy in the language divided by accuracy in English, clipped below at 0.05, worst measured value for unlisted languages, flagged), its distribution across countries is reported in \nolinkurl{coverage.csv} before any run, and it enters H4 only as a pre-registered secondary variant $c^{+}_k = (q^{\text{demo}}_k)^{0.25}(q^{\text{att}}_k)^{0.5}(q^{\text{lang}}_k)^{0.25}$ and in H6, where native-language prompting would use it.

An earlier draft also included a component for the share of items with an independent validation target. It is removed from $c_k$, because it measures whether a country can be checked rather than how likely it is to be wrong; it is reported as a separate column. The weights, $\tau$ and thresholds are fixed here and not tuned on evaluation data. The score is a pre-registered heuristic for where the population is likely to be wrong, not a probability; H4 tests whether it predicts error, and if it fails the release drops it and keeps the tiers.

\paragraph{What we can say before building.} GLOPOP-S has own-survey demographics for 78 countries \citep{ton2024glopop} and WVS/EVS reach 115 to 120 countries or territories over four decades \citep{haerpfer2022wvs}, but how many countries fall into each tier is an output of the build. It is computed from metadata alone, before any outcome is seen, and committed in \nolinkurl{artifacts/global-population/v0/coverage.csv}, which also fixes how many units each hypothesis can evaluate.

\subsection{Release and disclosure control}
\label{sec:release}

An agent row holding country, region, sex, age, education and a donor's full WVS item vector \emph{is} the donor's record: anyone with the public file can match the vector back to the respondent. Releasing agent-level attitude vectors would therefore redistribute survey microdata, which the WVS conditions forbid \citep{wvsconditions}, and would defeat any promise to withhold donor identifiers. We will not release agent-level attitude vectors from any source.

The public release contains: the code; the schema; the fallback map; a datasheet per country (Appendix~\ref{app:datasheet}) with source, vintage, $n_k$, $m_k$, tier, confidence components, total standard errors (Section~\ref{sec:uncertainty}) and Fréchet bound widths; agent-level demographic rows only if counsel confirms that GLOPOP-S may be redistributed in derived form given its upstream terms; a recipe that rebuilds the attitude block from survey files the user downloads under their own license; and aggregate attitude tables over cells of at least ten respondents, only where the source's terms permit publishing derived statistics for our use.

Inside the product, an agent's explanation shows provenance (source, wave, match key) and attitude items only from a fixed \emph{display set}: at most $B = 5$ items per agent, drawn once at build time from that agent's general-referent items and never changed, so no sequence of queries reveals more than those $B$ items. Any displayed or released format $F$ passes a pre-registered disclosure test before use, a $k$-anonymity check \citep{sweeney2002} with the same threshold as aggregate tables: for every record in $F$, count the respondents in the public donor file who share all of $F$'s fields (quasi-identifiers plus the full display set); $F$ is used only if every count is at least ten. Agents that fail have their display set shrunk until they pass. Results go to \nolinkurl{artifacts/global-population/v0/disclosure.json}.

\section{Cultural flattening: threat model and diagnostics}
\label{sec:flattening}

We call a population \emph{flattened} when its countries are more alike, or its people more alike within a country, than real surveys show, and \emph{caricatured} when countries are pushed further from the world mean than the data show (prompted country shifts can reflect stereotypes \citep{durmus2023}). If matching works as designed, the donor-carried part of the attitude block inherits the diversity of the donor survey, up to resampling roughly 800 donors with heavy cell collapse (Section~\ref{sec:stage2}); whether it does is measured by H2. Flattening can return in three places: borrowed donors (a country inherits its neighbour's distribution), the LLM's answers to items not in the donor vector (where the model's default takes over), and the daily belief updates (outside this paper). We measure the first two.

For item $j$, let $\mu^{\text{sim}}_{k,j}$ and $\mu^{\text{hum}}_{k,j}$ be the simulated and human country means (survey-weighted for humans, design-weighted for agents), with standard errors $s^{\text{sim}}_{k,j}$ and $s^{\text{hum}}_{k,j}$, and $\bar\mu_j$ the population-weighted mean of $\mu^{\text{hum}}_{k,j}$ over evaluated countries. Write $d^{\text{hum}}_{k,j} = \mu^{\text{hum}}_{k,j} - \bar\mu_j$ and $d^{\text{sim}}_{k,j} = \mu^{\text{sim}}_{k,j} - \bar\mu_j$.

\paragraph{Between-country slope.} $\beta_j$ is the slope of $d^{\text{sim}}_{k,j}$ on $d^{\text{hum}}_{k,j}$ across countries, divided by the reliability $\lambda_j = 1 - \overline{(s^{\text{hum}}_{k,j})^2}/\operatorname{Var}_k(d^{\text{hum}}_{k,j})$ to undo attenuation from human sampling error (which would otherwise mimic flattening). $\beta<1$ is flattening; $\beta>1$ is caricature. The summary $\beta$ is the median over items. Items with $\lambda_j < 0.5$ are excluded from every slope and ratio summary, and their number is reported, so near-singular corrections cannot drive the median. It replaces a per-country ratio $|d^{\text{sim}}|/|d^{\text{hum}}|$, which is unstable for countries near the world mean; per-country caricature is reported only as the signed residual for countries with $|d^{\text{hum}}_{k,j}| > 2 s^{\text{hum}}_{k,j}$.

\paragraph{Noise-corrected variance ratios.}
\begin{equation}
\rho^{B}_j = \frac{\operatorname{Var}_k\big(\mu^{\text{sim}}_{k,j}\big) - \overline{(s^{\text{sim}}_{k,j})^2}}{\operatorname{Var}_k\big(\mu^{\text{hum}}_{k,j}\big) - \overline{(s^{\text{hum}}_{k,j})^2}},\qquad
\rho^{W}_{k,j} = \frac{\operatorname{Var}\big(Y^{\text{sim}}_{j}\mid k\big)}{\operatorname{Var}\big(Y^{\text{hum}}_{j}\mid k\big)},
\end{equation}
on ordinal items scored in $[0,1]$. $\rho^B<1$ is between-country flattening and $\rho^W<1$ within-country variance collapse. The $\lambda_j \ge 0.5$ rule implies a positive $\rho^B$ denominator of at least half the observed between-country variance. Both ratios are estimated without bias correction of the ratio itself and carry cluster-bootstrap intervals.

\paragraph{Default-ward shift.} With JS the Jensen--Shannon distance and $p^{0}_j$ the model's own answer distribution to item $j$ with no persona and no country,
\begin{equation}
\delta_{k,j} = \widetilde{\operatorname{JS}}\big(p^{\text{sim}}_{k,j},\,p^{0}_j\big) - \widetilde{\operatorname{JS}}\big(p^{\text{hum}}_{k,j},\,p^{0}_j\big),
\end{equation}
where $\widetilde{\operatorname{JS}}$ is JS net of its noise ceiling at that pair's own sample sizes (Section~\ref{sec:noise}), so the two terms are comparable although they rest on different sample sizes. $\delta<0$ means simulated answers sit closer to the model's own default than the country's respondents do. This reference needs no assumption about which population the model favours. As a secondary reference we use US respondents, one of the populations closest to default model answers in \citet{durmus2023}; with that reference, $\delta<0$ can also reflect real convergence of a country toward the US, so it is not read as model drift. $\delta$ carries a cluster-bootstrap interval.

\section{Pre-registered evaluation protocol}
\label{sec:protocol}

\subsection{Registry: what is fixed, and when}
\label{sec:registry}

Table~\ref{tab:registry} lists every analytic choice. Elements marked \emph{here} are fixed by this paper. Elements marked \emph{registry} are written to \nolinkurl{artifacts/global-population/v0/registry.json} before the first evaluation run; that file's SHA-256 hash is committed to the public repository and posted publicly, and the evaluation harness refuses to start if the file's hash does not match. Registry elements depend only on metadata (questionnaires, file lists, release dates), never on outcomes. For each registry element this paper fixes the selection \emph{rule}; the resulting content (for example which items enter $J_c$) cannot be judged at Stage~1, and in-principle acceptance covers the rules, not that content.

\paragraph{Deviations.} If any element must change after Stage-1 acceptance, the change is written as a versioned amendment to the registry, with its reason and date, committed and hashed before any outcome it could affect is unblinded, and listed in the Stage-2 report next to the original.

\begin{table}[!htbp]
\centering
\small
\caption{Frozen elements. ``Here'' = fixed in this paper; ``registry'' = content frozen and hashed before any evaluation run under the rule given here.}
\label{tab:registry}
\begin{tabular}{@{}L{5.0cm}L{1.4cm}L{8.0cm}@{}}
\toprule
Element & Fixed & Where / rule \\
\midrule
Inclusion rule, $N$, $f$, frame & here & Section~\ref{sec:alloc} \\
Eligible wave, donor wave, recent wave & here & Section~\ref{sec:stage2} \\
Tier decision list, score form, weights, $\tau$, $b$ & here & Section~\ref{sec:tiers} \\
Fallback rule and alternative rule & here & Section~\ref{sec:fallback}; maps in \nolinkurl{fallback_map.csv} \\
Education crosswalk (GLOPOP-S to ISCED 2011) & registry & set-valued table, Section~\ref{sec:stage2} \\
H1 census cell crosswalk & registry & fixed cell sets, H1 \\
Common item set $J_c$ and cross-program crosswalk & registry & rules (i)--(iii), Section~\ref{sec:items} \\
Referent classes of items & registry & two coders, stricter class wins, Section~\ref{sec:items} \\
Held-out item rule (within-wave) & here & random 25\% of the donor wave's general items, stratified by topic; seed in registry \\
Post-cutoff target list, data cutoff date $T_D$ & here / registry & rules in Section~\ref{sec:splits}; list in registry \\
Minimum units per hypothesis, downgrade rule & here & Table~\ref{tab:minima} \\
Primary model & here & Claude Haiku 5.5, training data cutoff June 2026 \citep{claudemodels}; snapshot id in registry \\
Replication model & registry & the most recent open-weights instruction-tuned model with a stated training cutoff, named before the run \\
Prompt templates, elicitation settings & here / registry & settings in Section~\ref{sec:elicit}; templates in registry \\
$q^{\text{lang}}$ item set & registry & 200 MMMLU items, Section~\ref{sec:tiers} \\
H6 country list (20) & registry & spans tiers, scripts and $q^{\text{lang}}$ range \\
H1 displacement list & here & rule in H1; list in registry \\
Statistical tests, $\alpha$, families, SESOIs & here & Tables~\ref{tab:tests} and~\ref{tab:minima} \\
Compute budget & registry & Section~\ref{sec:power} \\
Random seeds & registry & all splits, draws and bootstraps \\
Power simulation & registry & \nolinkurl{power.json}, run after \nolinkurl{coverage.csv}, before any evaluation \\
\bottomrule
\end{tabular}
\end{table}

\subsection{Eliciting answers}
\label{sec:elicit}

Each evaluation agent answers each item once per donor draw, so each agent gives $M = 5$ answers per item. The prompt states the agent's profile and the item with its full response list, in an option order randomized per agent and item with a seeded generator, because option order and labels move model survey answers \citep{dominguez2024,tjuatja2024}; the don't-know option is included where the questionnaire offers it. Claude Haiku 5.5 has adaptive thinking \citep{claudemodels}; the thinking and effort settings are fixed in the registry, temperature is 1.0 where the API accepts it in that mode, and if it does not, the setting actually used is recorded in the registry before any run. We read the sampled answer, not token probabilities, so the same readout works for every model. An unparseable answer is retried up to twice with the same prompt; a still-invalid answer, or a refusal, is recorded as such (separately) and counted. In the primary analysis these are excluded from distributions, and a country whose invalid-plus-refusal rate on an item set exceeds 5\% is reported as not evaluable for that set; a registered sensitivity analysis keeps those countries and scores refusals as their own category, and every hypothesis lists the countries excluded for this reason. Baseline B3 asks the model directly for the country's answer distribution as stated probabilities per option, with the same option randomization. Every condition is run on the primary model; H3--H5 are repeated on the replication model.

\subsection{Splits, targets, item sets and feasibility}
\label{sec:splits}

\paragraph{Within-wave split.} For every country with an own donor wave, from whichever program supplies it, respondents are split 80/20 by primary sampling unit: whole PSUs go to one side, stratified by region, so donors and held-out humans never share a cluster. Where a file does not identify PSUs the split is by respondent within region and $X$ cells; such countries are flagged, because intra-cluster correlation then flatters H2 and set W, and a sensitivity analysis drops them. Donors come only from the 80\%; the 20\% are held-out humans.

\paragraph{Evaluation agents.} For each target sample, $n^{\text{eval}} = \min(n^{\text{tgt}}, 300)$ target respondents are drawn by weighted systematic sampling, and an evaluation agent is built from each one's $X$, so the comparison is between populations with the same demographic make-up.

\paragraph{Item sets.} The primary model's training data cutoff is June 2026 \citep{claudemodels}; we cannot verify it, and effective cutoffs can differ from stated ones \citep{cheng2024dated}.
\begin{itemize}
  \item[(P)] \emph{Post-cutoff targets}: target waves whose fieldwork began after the cutoff and whose microdata are released and licensed by the data cutoff date $T_D$. $P$ has two strata. \emph{$P_{\text{rep}}$}: general items repeated from earlier waves (in $J_c$ or a program's common items), masked from the agent's donor vector, so neither agent nor baseline sees a donor answer to the item, and the target answers postdate the model's training. \emph{$P_{\text{new}}$}: items whose wording appears in no questionnaire published before the cutoff. $P_{\text{new}}$ is the product's novel-question situation but rarely has cross-country coverage; $P_{\text{rep}}$ supports cross-country statistics. Older answers to $P_{\text{rep}}$ items may be memorized; that knowledge is also available to deployable baselines and to the deployed system, so it does not bias comparisons among arms.
  \item[(I)] \emph{Independent pre-cutoff targets}: the same kind of masked repeated items scored against an independent wave fielded before the cutoff (a newer wave of the same or another program). Contamination is possible.
  \item[(W)] \emph{Within-wave held-out items}: the held-out 25\% of the donor wave's general items, masked from donor vectors and scored against the 20\% holdout. Its items are public and may be memorized \citep{durmus2023,zhao2024wvb}.
\end{itemize}

\paragraph{Feasibility of set P.} With a June 2026 cutoff, only waves fielded from mid-2026 whose microdata are released before $T_D$ qualify. We expect $P$ to be small and it may be empty: for example, Arab Barometer Wave IX was fielded from September 2025 to May 2026 \citep{arabbarometer}, before the cutoff, and contributes nothing. The primary model is to be served until at least 7 October 2027 \citep{claudemodels}; we set $T_D$ = 7 July 2027, three months before that date. Before any run, \nolinkurl{coverage.csv} lists every program wave with fieldwork after the cutoff, its countries, its announced or expected release date where the program states one, and the resulting candidate counts per hypothesis. Table~\ref{tab:minima} gives, for each hypothesis that uses $P$, the minimum units required for its primary test to be confirmatory. If a minimum is not met at $T_D$, that hypothesis is run on its fallback set and \emph{labelled exploratory} in the Stage-2 report; this is a pre-registered downgrade, not a switch of primary set. A confirmatory test can then be re-registered, by amendment, for a successor model with a later cutoff.

\begin{table}[!htbp]
\centering
\small
\caption{Registered minima for confirmatory primary tests, smallest effects of interest (SESOI) used by the power simulation, and the fallback when a minimum is not met at $T_D$.}
\label{tab:minima}
\begin{tabular}{@{}L{0.7cm}L{4.2cm}L{4.6cm}L{4.4cm}@{}}
\toprule
 & Minimum for confirmatory & SESOI & If not met \\
\midrule
H1 & 15 countries with a census in or after 2015 & 5\% median relative SRMSE reduction & exploratory, all available censuses \\
H2 & 20 countries with $\ge 2$ hidden variables & share of severely attenuated pairs 0.35 vs.\ tolerance 0.5 & exploratory \\
H3 & 20 countries with $\ge 10$ items in $P$ & 10\% relative reduction in mean $e$ vs.\ B3 in the median country & exploratory on I, then W \\
H4 & 55 countries with $\ge 10$ $P_{\text{rep}}$ items & partial $\rho = -0.45$ & exploratory on I \\
H5 & 15 countries sharing $\ge 10$ $P_{\text{rep}}$ items & $\Delta\beta = 0.15$ & exploratory on I \\
H6 & 10 of the 20 countries with $\ge 10$ items in $P$ & 10\% relative change in mean $e$ & exploratory on I \\
H7 & 10 countries with $\ge 3$ qualifying waves & slope of 2\% of $e_0$ per year & exploratory \\
H8 & 20 countries with $\ge 5$ mapped topline items & non-inferiority margin 2 points & exploratory \\
\bottomrule
\end{tabular}
\end{table}

\paragraph{Simulated borrowing.} For every country with an own donor wave and an eligible fallback, we also build the country as if it had no own data, with donors from $k'$ under the committed fallback rule, and score it against the country's own targets. This gives Tier-3-like units with real targets, a direct estimate of the borrowing penalty, and a measurement of flattening through borrowing. Its confidence score uses $b_k=0.5$ and the fallback wave's $t_{k'}$ and $m_{k'}$. The arm is built only from data-rich countries whose neighbours also have data, while real Tier-3 countries tend to be smaller, poorer, more often conflict-affected and to have sparser neighbours; the measured borrowing penalty is therefore likely a lower bound. The arm evaluates borrowed attitudes only, not the filled or WPP-only demographics that also place countries in Tier~3.

\subsection{Error metric and noise ceiling}
\label{sec:noise}

For a country, item and arm, raw error is the JS distance between the arm's answer distribution and the target's. Plug-in JS between finite samples is biased upward by an amount that depends on both sample sizes and on the number of categories; because target size tracks donor-wave size, raw error would also correlate with the score's sample-size factor for reasons unrelated to quality. Every error is therefore reported net of a \emph{noise ceiling} computed at the comparison's own sample structure: $\operatorname{JS}_{\text{ceil}}$ is the mean, over 200 resamples, of the JS distance between two weighted resamples drawn from the target itself, one of size $n^{\text{tgt}}$ and one with the arm's structure. For agent arms the second has $n^{\text{eval}}$ units; we treat each agent's five correlated draws as carrying one answer's information, which makes the ceiling larger and the corrected error smaller, and report a sensitivity analysis at $5n^{\text{eval}}$. For B3, which states probabilities, the arm side has no sampling noise. For B1 and B2 the second resample has the oracle's own size. The outcome is noise-corrected error $e = \operatorname{JS} - \operatorname{JS}_{\text{ceil}}$; negative values are kept, not truncated. The same correction gives $\widetilde{\operatorname{JS}}$ in $\delta$.

Paired comparisons between arms on the same items (H3, H5, H6) need no further normalisation. Where errors are compared across countries whose item sets differ (H4, H7), the outcome is item-centred error, $e_{k,j} - \bar e_j$ with $\bar e_j$ the mean over countries for item $j$, averaged per country, so items with more categories or greater difficulty do not drive country differences. As a registered sensitivity analysis for ordinal items, every error is recomputed with the Wasserstein-1 distance on the $[0,1]$ scores, which respects order. All human distributions, B1 and the ceiling use survey weights; agent distributions use design weights.

\subsection{Uncertainty}
\label{sec:uncertainty}

\paragraph{Released estimates.} For a country estimate $Q$ (a proportion or mean of a donor-carried item), each of the $M=20$ proper donor draws (Section~\ref{sec:stage2}) gives $\hat Q_m$ and an agent-sampling variance $U_m$ from the design weights. Rubin's rules \citep{rubin1987} give $\bar Q = M^{-1}\sum_m \hat Q_m$, $\bar U = M^{-1}\sum_m U_m$, $B = (M-1)^{-1}\sum_m(\hat Q_m-\bar Q)^2$ and total variance $T = \bar U + (1+M^{-1})B$. Because each draw resamples the donor wave by PSU or replicate weights, $B$ carries the donor survey's sampling and design error, which dominates $\bar U$ for any country with $n_k \gg m_k$. The release ships $\sqrt{T}$ as the \emph{total standard error} of every country estimate; it does not ship a ``Monte Carlo'' error that omits the donor survey.

\paragraph{Evaluation.} Intervals come from a cluster bootstrap with 2,000 replicates that jointly resamples target respondents within PSUs where identified, items, and donor draws. Country-level tests (H2--H5, H7, H8) treat the country as the unit: their inference resamples or permutes countries, as stated per test. A mixed-effects model with country and item random effects is reported as a robustness check.

\subsection{Baselines}

Two baselines are \emph{oracle references} that use answers to the target item and so are unavailable in deployment: (B1) the weighted country marginal of the item from the donor 80\%; (B2) the matched donor's own answer to the item. They bound what is achievable on set W and are not competitors. Four baselines are \emph{deployable}, never seeing answers to the target item in country $k$: (B3) the model asked directly for the country's answer distribution \citep{jang2026,meister2025}; (B4) a demographics-only persona with no donor answers \citep{argyle2023}; (B5) the nearest-item rule: the weighted marginal, in $k$'s donor wave, of the conditioning item most similar to the target item under a committed text-embedding similarity, mapped onto the target scale by matching cumulative shares; and (B6) a leave-one-country-out multilevel regression with poststratification \citep{park2004mrp,ghitza2013}, fitted on other countries' respondents to the same item with $X$ and conditioning items as predictors and poststratified to $k$'s evaluation $X$ (available only where other countries answered the item).

\subsection{Hypotheses}

Each hypothesis has one primary test (Table~\ref{tab:tests}); other predictions are secondary and reported without entering the error-controlled family.

\begin{description}
  \item[H1 Skeleton fidelity: does cohort shift help?] For countries with a census in or after 2015 that publishes the needed tables (or IPUMS International aggregates if permission is granted), plain and cohort-shift skeletons are both post-stratified to WPP totals \emph{for the census year}, so the comparison isolates the method from change since the census. A committed crosswalk maps every census table to two fixed cell sets: 6 age bands $\times$ 2 sexes $\times$ 3 education groups (GLOPOP-S levels 1--2, 3, and 4--5, mapped to ISCED 2011 levels 0--1, 2 and incomplete 3, and complete 3 or above \citep{isced2011}), and 6 age bands $\times$ 2 sexes $\times$ 2 settlement types, so every country has $|C| = 36$ and $|C| = 24$ cells and SRMSE $= \sqrt{|C|^{-1}\sum_c(\hat p_c - p_c)^2}\,\big/\,\big(|C|^{-1}\sum_c p_c\big)$ is comparable across countries. Primary: the paired per-country difference in SRMSE (cohort shift minus plain, averaged over the two tables). Decision rule: a new, versioned release switches to cohort shift if the one-sided Wilcoxon signed-rank test (cohort shift lower) rejects at the Holm-adjusted level and the median paired relative reduction is at least 5\%; otherwise plain reweighting stays. Secondary: error for GLOPOP-S fill versus own-survey countries (descriptive, since census vintage and availability confound it); plain-reweighting error for pre-2015 censuses; error growth with years since 2015 among countries with large displacement, defined as absolute cumulative WPP net migration between 2015 and the census year of at least 5\% of the 2015 population \citep{wpp2024}. \pending{artifacts/global-population/v0/eval/h1_skeleton.json}.
  \item[H2 Matching and the conditional-independence assumption.] Within the donor wave, agents built from holdout $X$ with donors from the 80\%. Matching on $X$ can only carry the part of a hidden variable's association with attitudes that runs through $X$, so some attenuation is guaranteed; the question is how much. Hidden variables are survey variables deliberately left out of the key: GDL region (where the survey's region variable maps to it), income or wealth scale, marital status, employment status and household size, where the wave carries them. For each item--hidden-variable pair whose holdout bias-corrected Cramér's $V$ \citep{bergsma2013} is at least 0.10, the attenuation ratio is bias-corrected $V$ in the matched file over that in the holdout. Primary: the share $S$ of such pairs, pooled across countries with each country weighted equally, whose ratio is below 0.5 (``severely attenuated''). One-sided test of $H_0\!: S \ge 0.5$ against $S < 0.5$ by cluster bootstrap over countries; a pass means matching keeps at least half the strength of most hidden associations. Decision rule: within each country, a hidden variable whose median ratio is below 0.5 is labelled ``matching-assumption dominated'', and product cross-tabulations by that variable (including region) carry that label. Secondary (H2b): at the $X$-cell level (cells with at least 20 holdout respondents), matched agents beat B1 on noise-corrected JS; at the country level matching is expected to tie B1 by construction, and is reported only as a check. This tests the matching, not the LLM. \pending{artifacts/global-population/v0/eval/h2_matching.json}.
  \item[H3 Agent answers on items with no donor answer.] Donor-grounded agents answer set $P$ items (both strata, reported separately); noise-corrected error is compared with B3--B6 per country, and B1--B2 are reported on set W as upper-bound references. Primary: across countries, the per-country point difference in mean noise-corrected error (grounded minus B3) is negative in more than half of countries; one-sided sign test on the signs of these point differences, null probability 0.5. Separately, per-country wins form their own family: a country is a win if its one-sided p-value for a negative difference, from a respondent-cluster bootstrap, survives Benjamini--Hochberg control at $q = 0.10$ across countries. Secondary: the same against B4, B5, B6; a contamination probe comparing B3's error on the same $P_{\text{rep}}$ items scored against pre-cutoff (I) and post-cutoff (P) targets, which holds item difficulty and topic fixed but not time, so it is only suggestive of memorization. Decision rule: per country, the better of B5 and B6 is chosen on set I (or W where I is missing), never on $P$; where grounded agents do not beat that baseline on $P$ with a per-country BH-controlled win (a second per-country family), downstream products use that baseline for novel items in that country. \pending{artifacts/global-population/v0/eval/h3_heldout.json}.
  \item[H4 Confidence predicts error.] Units are real builds of Tier 1--3 countries, one unit per country, each with one outcome: item-centred noise-corrected error on its $P_{\text{rep}}$ items against post-cutoff targets, so staleness and borrowing can both affect error. $c_k$ is recomputed with $t$ at the target's fieldwork midpoint. Primary: Spearman partial correlation between $c_k$ and error across countries, controlling for $\log n^{\text{tgt}}$. Inference: Freedman--Lane permutation \citep{freedman1983} at country level (rank-transform, regress ranked error on ranked $\log n^{\text{tgt}}$, permute residuals across countries, refit; 10,000 permutations). Decision rule: the score is published only if the one-sided upper confidence bound at the Holm-adjusted level is below zero \emph{and} the point estimate is $\le -0.3$; otherwise it is reported as failed and removed. Because $c_k$ contains $\min(1, m_k/1000)$ and target and donor often come from one program, controlling for target size may remove real signal; the Spearman correlation between $\log n^{\text{tgt}}$ and $\log m_k$ is reported in \nolinkurl{coverage.csv} before the run, and the unadjusted correlation is reported beside the primary. Secondary: a two-level analysis pooling real and simulated-borrowing builds (country random intercept; between-country and within-country slopes reported separately, the within-country part being the borrowing contrast); the borrowing penalty (paired real minus simulated-borrowing error); tier ordering (Tier 1 $<$ 2 $<$ 3 error, Jonckheere--Terpstra); the variant $c^{+}_k$ with $q^{\text{lang}}$; partial correlations of single components ($m_k$, years since wave, borrowing) to show whether the composite adds anything; and rank stability (Kendall's $\tau$) of $c_k$ under $\pm50\%$ perturbations of each weight and $\tau \in \{4, 16\}$, not used for tuning. \pending{artifacts/global-population/v0/eval/h4_confidence.json}.
  \item[H5 Flattening.] On $P_{\text{rep}}$ items, with evaluation agents under grounding and under B4. Primary: $\Delta\beta = \beta_{\text{grounded}} - \beta_{\text{B4}} > 0$ (grounding brings the corrected slope closer to one), tested with a country-level paired bootstrap: countries are resampled with replacement and both slopes recomputed on the same resampled countries, 2,000 replicates, one-sided percentile interval. Items with $\lambda_j < 0.5$ are excluded (Section~\ref{sec:flattening}). Secondary: B4 shows $\beta<1$, $\rho^B<1$ and $\delta<0$, consistent with \citet{durmus2023,tao2024}; grounding moves $\log\rho^B$ and $\delta$ toward zero; $\beta$ and $\delta$ are worse in simulated-borrowing builds than in the same countries' real builds (flattening through borrowing); $\rho^W<1$ for both, as \citet{jang2026} predicts. \pending{artifacts/global-population/v0/eval/h5_flattening.json}.
  \item[H6 Native-language agents.] Twenty countries from the registry, in a $2\times2$ design: prompt language (English, native) $\times$ framing (persona only, persona with explicit cultural framing). Native-language prompts use the program's own national-language questionnaire wording where the program publishes one; items that exist only in machine translation are excluded from the primary analysis and analysed as a separate stratum, so translation quality does not masquerade as a language effect. Outcomes: noise-corrected error on set $P$ (primary), $\rho^W$, and measured input and output tokens per answer. Primary: the pooled main effect of native language within the framed condition (native-with-framing minus English-with-framing), two-sided, paired bootstrap over countries, with the language $\times$ framing interaction reported. Predictions from the literature: framing $\geq$ native language on average \citep{bulte2026,kwok2024}, with country exceptions \citep{alkhamissi2024}. Whether native language raises within-country variance is exploratory. Decision rule: native language is enabled for a country only if native-with-framing beats English-with-framing with a per-country one-sided p-value, from a respondent-cluster bootstrap, that survives Benjamini--Hochberg control at $q = 0.10$ across the 20 countries. \pending{artifacts/global-population/v0/eval/h6_language.json}.
  \item[H7 Staleness.] Countries with at least three eligible waves of the same program and the same interview mode. The item set is fixed per country: the general items present in every wave used for that country. Each older wave serves as the donor wave and is scored against the newest. Primary items are \emph{donor-carried}: agents copy the donor answer, as released agents do. H7 therefore measures how fast copied answers go stale, which is human attitude change filtered through the skeleton and matching; it is the quantity the recency factor of $c_k$ models, and it does not test the LLM. Primary: in a mixed model of item-centred noise-corrected error on the gap in years with a country random intercept, the gap slope is positive (one-sided). Secondary: the same with the items masked and answered by the model; and a fit of $e(\text{gap}) = e_\infty - (e_\infty - e_0)\,e^{-\text{gap}/\tau}$ by nonlinear least squares. The fitted $\tau$ equals the score's $\tau$ only under the assumption that the quality the score tracks is proportional to $(e_\infty - e)/(e_\infty - e_0)$; we state that assumption rather than test it. A ``mismatch'' with $\tau = 8$ is declared if the 95\% interval for $\tau$ excludes 8, and leads to a new, versioned score, not a silent retune. With three to five waves and gaps of at most 15 years the interval will be wide, and failing to exclude 8 is not evidence that $\tau = 8$. Waves with a mode change are analysed separately. A cohort-aware variant that matches on birth-cohort band (age at the target minus the gap) instead of age band is reported as secondary, since a 30-year-old in 2026 matched to a 2011 wave's 30-year-old belongs to a different cohort. \pending{artifacts/global-population/v0/eval/h7_staleness.json}.
  \item[H8 The released object.] H2--H7 use evaluation agents built from survey respondents' $X$, not the 100,000 released agents. H8 tests the release itself against published Gallup or Pew country toplines we may use. Items are topline items that map one-to-one onto a donor-carried general item; the release estimate is the design-weighted aggregate over all agents and $M=20$ draws, needing no model calls. Frames are aligned: a topline is compared only with agents in the same age frame (for Gallup, 15 and over \citep{gallupwp}); where the agent frame is 18 and over and the topline is 15 and over, the pair is used only if the topline publishes an 18-and-over figure, and otherwise reported in a frame-mismatch stratum. Primary: non-inferiority of the build relative to its own donor survey. For each country--item pair, $\Delta = |\hat Q^{\text{release}} - Q^{\text{topline}}| - |\hat Q^{\text{donor}} - Q^{\text{topline}}|$ in percentage points, with $\hat Q^{\text{donor}}$ the donor wave's own survey-weighted estimate. Survey house, mode and wording differences between the donor program and the topline enter both terms, so $\Delta$ isolates what the build adds (frame alignment, reweighting to the 2026 skeleton, matching and collapse). $H_0\!: \bar\Delta \ge 2$ against $H_1\!: \bar\Delta < 2$, with $\bar\Delta$ the mean over countries of per-country means, one-sided, cluster bootstrap over countries. Secondary: coverage of the topline by the release interval widened by an inter-survey discrepancy variance estimated, leaving the country out, from the donor-versus-topline differences of the same items in other countries, with nominal 95\% coverage as the target; and a model-answered stratum, topline items with no donor counterpart answered by a fixed random subsample of 100 agents per country with five draws each, reported and not confirmatory. \pending{artifacts/global-population/v0/eval/h8_release.json}.
\end{description}

\subsection{Error control and decision rules}

Table~\ref{tab:tests} gives one primary test per hypothesis. The eight primary tests form one family controlled by Holm's procedure \citep{holm1979} at familywise $\alpha = 0.05$. Holm across all eight is valid; a primary that rejects easily loosens the thresholds for the rest, which is one more reason the primaries are magnitude or paired tests rather than tests the design guarantees. Per-country decisions (H3 wins over B3, H3 wins over the better non-LLM baseline, H6 enablement) are three separate families, each controlled by the Benjamini--Hochberg procedure \citep{benjamini1995} at $q = 0.10$ applied to per-country p-values. For H6 this bounds the expected share of falsely enabled countries among the countries enabled at 10\%; it does not bound how many countries that would benefit are missed. Secondary results are reported with unadjusted intervals and labelled as such. A hypothesis that misses its minimum in Table~\ref{tab:minima} leaves the Holm family and is reported as exploratory; Holm is then applied to the remaining primaries.

\begin{table}[!htbp]
\centering
\small
\caption{Primary statistic and test for each hypothesis. Unit of inference is the country unless stated.}
\label{tab:tests}
\begin{tabular}{@{}L{0.7cm}L{4.9cm}L{4.6cm}L{1.0cm}L{2.7cm}@{}}
\toprule
 & Primary statistic & Test & Sides & Family \\
\midrule
H1 & Paired SRMSE, cohort shift minus plain & Wilcoxon signed-rank; switch if rejects and median relative reduction $\ge 5\%$ & one & Holm, $\alpha=0.05$ \\
H2 & Share $S$ of pairs with bias-corrected $V$ ratio $< 0.5$ & Cluster bootstrap over countries, $H_0\!: S \ge 0.5$ & one & Holm \\
H3 & Signs of per-country point differences, grounded minus B3 & Sign test, null 0.5 & one & Holm; per-country wins BH $q=0.10$ \\
H4 & Partial Spearman($c_k$, $e$ $\mid$ $\log n^{\text{tgt}}$), real builds & Freedman--Lane permutation; keep if upper bound $<0$ and estimate $\le -0.3$ & one & Holm \\
H5 & $\beta_{\text{grounded}} - \beta_{\text{B4}}$ & Country-level paired bootstrap & one & Holm \\
H6 & Native minus English, framed arms & Paired bootstrap over countries & two & Holm; per-country BH $q=0.10$ \\
H7 & Gap slope in mixed model, donor-carried items & Likelihood-ratio / Wald & one & Holm \\
H8 & Mean $\Delta$, release minus donor-survey absolute error & Non-inferiority, margin 2 points, cluster bootstrap & one & Holm \\
\bottomrule
\end{tabular}
\end{table}

\subsection{Power, effect sizes and budget}
\label{sec:power}

The number of evaluable units for each hypothesis is fixed by \nolinkurl{coverage.csv} before any outcome is seen. Power needs an assumed alternative, so the smallest effect of interest for each hypothesis is fixed in Table~\ref{tab:minima}, and the generative model for the simulation is fixed here: noise-corrected error $e_{k,j,a} = \mu + u_k + v_j + \gamma_a + \eta\,c_k + \varepsilon_{k,j,a}$, with country, item and arm effects, the variance components $u_k$, $v_j$ and $\varepsilon$ estimated from split-half resamples of human data only (which contain no model output), and the arm and score effects $\gamma_a$, $\eta$ set to the SESOIs. The simulation estimates, for each hypothesis, power at its SESOI given the actual units and holdout sizes. Any hypothesis whose simulated power falls below 0.8 is relabelled exploratory in the registry before the run. The simulation is committed as \nolinkurl{artifacts/global-population/v0/power.json}.

As design arithmetic for H4, with a Fisher-$z$ approximation for Spearman's partial correlation using variance $1.06/(n-4)$, the decision rule (upper bound below zero at one-sided $\alpha = 0.05/8$, the strictest Holm level, \emph{and} point estimate $\le -0.3$) has 80\% power at a true partial correlation of $-0.45$ with about 55 countries, of $-0.5$ with about 44, and of $-0.4$ with about 70. At a true value of $-0.3$ the point-estimate condition alone caps power at 50\% for any number of countries. The units are countries with real builds; simulated-borrowing builds share countries and targets with them and do not add independent units to the primary.

\paragraph{Compute budget.} Model calls for evaluation are $\sum_{\text{arms}}\sum_k n^{\text{eval}}_k\,|J_k|\,M$ per model, over the per-agent arms (grounded, B4, the four H6 arms in 20 countries, simulated borrowing), with $|J_k|$ the evaluated items in country $k$ and $M = 5$, doubled for the replication model on H3--H5, plus one call per country and item for B3 and the unconditioned $p^0$ calls. The registry fixes a ceiling on calls and tokens from this formula. If the ceiling would be exceeded, items are removed in a random order fixed in the registry, the same order in every country, so a shortfall cannot shrink the item set selectively.

\begin{table}[!htbp]
\centering
\small
\caption{Pending results. No entry has been produced. Each will be filled only from the named artifact, at a named commit.}
\label{tab:pending}
\begin{tabular}{@{}L{5.0cm}L{8.6cm}@{}}
\toprule
Result & Future artifact \\
\midrule
Registry of frozen elements and its hash & \nolinkurl{artifacts/global-population/v0/registry.json} \\
Countries, $n_k$, $m_k$, tiers, score components, $q^{\text{lang}}$, $X$ distributions, collapse shares, set $P$ candidates, evaluable units & \nolinkurl{artifacts/global-population/v0/coverage.csv} \\
Fallback donor maps (committed pre-run) & \nolinkurl{artifacts/global-population/v0/fallback_map.csv} \\
License status per source and use & \nolinkurl{artifacts/global-population/v0/licenses.csv} \\
Power simulation & \nolinkurl{artifacts/global-population/v0/power.json} \\
Fréchet bound widths & \nolinkurl{artifacts/global-population/v0/frechet.csv} \\
Disclosure test & \nolinkurl{artifacts/global-population/v0/disclosure.json} \\
H1 cohort shift vs.\ plain skeleton & \nolinkurl{.../v0/eval/h1_skeleton.json} \\
H2 matching and hidden variables & \nolinkurl{.../v0/eval/h2_matching.json} \\
H3 items without donor answers vs.\ baselines & \nolinkurl{.../v0/eval/h3_heldout.json} \\
H4 confidence vs.\ error & \nolinkurl{.../v0/eval/h4_confidence.json} \\
H5 flattening diagnostics & \nolinkurl{.../v0/eval/h5_flattening.json} \\
H6 native-language arms, tokens & \nolinkurl{.../v0/eval/h6_language.json} \\
H7 staleness & \nolinkurl{.../v0/eval/h7_staleness.json} \\
H8 released object vs.\ donor surveys and toplines & \nolinkurl{.../v0/eval/h8_release.json} \\
\bottomrule
\end{tabular}
\end{table}

\paragraph{What would count against the design.} If grounded agents do not beat B3 on items with no donor answer, donor grounding adds nothing to the model's own estimate of a country, and the product should use direct distribution estimates. If they do not beat B5 or B6, the LLM adds nothing over a cheap statistical imputer. If H4 fails, per-country confidence cannot be computed from inputs, and the release says so. If H5 shows grounding does not reduce flattening, the ``whole world'' framing is not supported for the affected countries: the product describes those countries as simulated from model defaults, and their tier is what we report. If H8 fails, the build adds error to its own sources, and released country aggregates should be the donor surveys' estimates, not the population's.

\section{Relation to the rest of the system}
\label{sec:companions}

This paper covers who the agents are and how well each country is known. The system as a whole, belief updates, forecast commitment and scoring, dispersion correction of model answers, the benchmark, the review of the silicon-sample literature and the governance of public release are the subjects of separate papers in preparation (\companion{Don't Ask a Model, Ask a World}, \companion{Belief Assimilation}, \companion{Receipts}, \companion{Correcting Variance Collapse}, \companion{The Real World Benchmark}, \companion{Silicon Samples, 2022--2026}, \companion{Governing Public World Simulators}); nothing in this protocol depends on them, and the commitment procedure it uses is the one in Section~\ref{sec:registry}.

Three differences from the forecasting engine are deliberate, and the first is a gap. First, the engine never calls the model once per agent: its forecasting state is numeric, and the model is called per news story, per question or per segment (\companion{Don't Ask a Model, Ask a World}), as in archetype-based simulators \citep{chopra2024limits}. This paper evaluates per-agent answers. Whether results for per-agent answers carry over to segment-level calls is not tested here: the tiers, scores and provenance are properties of the population and can be attached to segments, but the error of a segment-level answer is a separate quantity that this protocol does not measure. Second, until the written permissions of Section~\ref{sec:sources} arrive, the engine's public build does not use donor-matched vectors at all; it is designed to anchor segments on published survey aggregates, and this resource replaces those anchors only for sources whose licence row reads \texttt{permitted}. The engine built so far does neither: its demographic cells, attitude shares and segment effects are illustrative values in \nolinkurl{data/sample/} that its licence table marks as a repository fixture, and its 16 countries get a coverage tier from a fixed five-part rubric score in \nolinkurl{data/countries.csv}; no population tier or confidence score of this paper is computed. Third, every product output carries a \emph{coverage tier} (A--D) defined by the attitude evidence behind the forecasting state; the population tiers here score only the population layer, and a Tier~4 country is coverage tier D. Native-language prompting is not part of the launch configuration (Section~\ref{sec:lang}).

\section{Limitations}
\label{sec:limits}

\paragraph{What is evaluated is not all of what ships.} H1--H8 test the population on survey items, through per-agent answers. The product mostly asks about novel events, updates beliefs daily, and queries the model per segment rather than per agent (Section~\ref{sec:companions}); until licences arrive it uses published aggregates rather than donor vectors. Set $P$ is the closest proxy we have, but passing it does not validate event questions, belief updates, segment-level answers or forecasts; those are evaluated by the forecast record, not here.

\paragraph{Set P may not exist in time.} With a June 2026 cutoff and a primary model served until at least October 2027 \citep{claudemodels}, few attitude waves may be fielded and released in the window. The downgrade rule (Table~\ref{tab:minima}) keeps the report honest, but the confirmatory evidence for H3--H6 may then wait for a successor model.

\paragraph{The matching assumption, and region.} Conditional independence given four demographic variables is almost certainly false for some attitude--attribute pairs. Region is the most consequential omission: in many countries region, ethnicity and religion carry much of the within-country variation in attitudes, and none is in the key. Any regional cross-tabulation in the product is, to the extent H2 labels region ``matching-assumption dominated'', an artifact of the assumption, and is labelled so. For skeleton attributes no survey carries, only Fréchet bounds apply.

\paragraph{Copied answers age and travel badly.} Referent classes remove the clearest invalid copies (a neighbour's view of its own government, a decade-old view of a former president), but classification is judgement, and general items can also change meaning across countries and time.

\paragraph{The skeleton is old and partly borrowed.} GLOPOP-S describes 2015 \citep{ton2024glopop}. Neither reweighting nor cohort shift recovers changes in urbanization or household structure, and both are weakest where change was fastest (conflict, displacement, large migration). Ninety-nine countries, 15.7\% of world population, are filled from three similar countries each \citep{ton2024glopop}. Education categories differ across sources, and the crosswalk can only make that ambiguity explicit; the cohort-shift education rule for the young is a stated approximation.

\paragraph{Borrowing is measured where it is easiest.} The simulated-borrowing arm uses data-rich countries with data-rich neighbours, so its borrowing penalty is likely a lower bound for real Tier-3 countries, and it does not evaluate filled demographics.

\paragraph{The surveys are imperfect ground truth.} National samples are around a thousand respondents, fielded by different modes and with different translations; some questions are sensitive or unaskable in some states, nonresponse is not random, and in authoritarian settings preference falsification can make a country look well measured when it is not. Agreement with a survey is agreement with that survey, not with the population. Excluding countries for refusals removes countries with sensitive items non-randomly; the sensitivity analysis in Section~\ref{sec:elicit} shows by how much.

\paragraph{Contamination.} Set $P$ guards against memorized answers only if the model's stated training cutoff is accurate; we cannot verify it, and effective cutoffs can differ from stated ones \citep{cheng2024dated}.

\paragraph{The score is a heuristic.} Its form and weights are judgement, fixed in advance to avoid tuning on evaluation data. A good H4 result would show it ranks countries usefully, not that it is calibrated.

\paragraph{Licensing.} Most sources restrict commercial use or redistribution (Section~\ref{sec:sources}). The protocol may not run at all for some sources, and the release may contain no agent-level rows.

\paragraph{The model moves.} Results apply to the model snapshot recorded in the run, and the primary model has a stated retirement date no earlier than 7 October 2027 \citep{claudemodels}. Any model change requires a rerun of H3--H6 before tiers or score are carried over; the replication model guards against results that hold for one vendor only.

\paragraph{Within-person consistency.} Each agent answers each item five times at temperature 1.0 under different donor draws; the evaluation measures distributions, not whether an agent answers consistently over time, which this design will tend to understate.

\paragraph{Cost.} The evaluation needs model calls on the order given by the budget formula in Section~\ref{sec:power}; its token cost will be reported from the run's logs, and the budget ceiling can remove items.

\paragraph{Scope.} Agents are persons aged 18 and over, or 15 and over where the attitude survey's frame starts at 15; they are not all adults. The resource makes no validity claim below GLOPOP-S regions. It represents a country as its survey-reachable population, which excludes many institutionalized, displaced and remote people.

\section{Ethics statement}

\paragraph{Conflict of interest.} The authors work for Real World, a commercial company that sells simulation outputs built on this population. We have a financial interest in H3 and H4 passing. The pre-registration, the registry hash, the deviation policy and the commitment to report every country and item, including failures, are the safeguards we offer. The primary model and the vendor table we cite \citep{anthropicmultilingual,claudemodels} are from the same vendor; we cite the table only for languages and do not use its numbers in the score.

\paragraph{Ethics review.} This is secondary use of de-identified survey files, but for a commercial purpose that respondents did not consent to. We have not yet obtained review. No donor file is ingested and no evaluation runs before an independent research ethics committee has approved the protocol; the committee, the submission date and its decision are recorded in the registry before it is frozen and reported at Stage~2. In-principle acceptance is sought conditional on that approval, and we state plainly that this is a gap at Stage~1.

\paragraph{Consent and custodian permission.} Survey respondents agreed to research uses described by each program, not to commercial simulation. A custodian's written permission makes our use lawful under the custodian's terms; it does not make it a use respondents consented to. We treat that gap as a reason for the release restrictions in Section~\ref{sec:release}, not as settled by permission. The same reasoning rules out using an academic collaborator's licence without the custodian's written confirmation (Section~\ref{sec:sources}).

\paragraph{Data protection.} Political opinions and religious beliefs are special categories of personal data under Article~9 of the EU General Data Protection Regulation \citep{gdpr}. EVS files and other European waves are de-identified, but donor vectors attached to agents inside a commercial system are still derived from such data. Before any EU-origin donor file is ingested we will document a lawful basis and complete a data protection impact assessment.

\paragraph{Respondent safety.} Respondents in authoritarian states answered sensitive questions under assurances of confidentiality. Copying their answer vectors into agents raises the cost of any re-identification. The release excludes agent-level attitude vectors, product displays show at most a fixed set of five items per agent, and the disclosure test gates every displayed format.

\paragraph{Dual use.} Country-specific populations that can be forked with an event are a ready tool for message testing and influence operations, most of all for low-tier countries where outputs are least checked. This paper commits to the following, which the release enforces: no agent-level attitude output leaves the system; every output carries its country's tier and score; Tier~4 countries are excluded from attitude outputs; displays pass the disclosure test; and no output is produced for a segment smaller than the disclosure threshold. Further product rules (for example on election-related outputs) belong to the governance paper and are not relied on here.

\paragraph{Speaking for populations.} A simulated country can be quoted as if it were the country. We will not present any agent output as a measurement of public opinion: published forecasts are forecasts of measurable outcomes, scored against those outcomes, and every population-derived figure is labelled simulated, with its tier and total standard error. If H5 fails for a country, the product describes that country's agents as reflecting model defaults rather than the country.

\paragraph{Representation and language.} Borrowing donors from a neighbour, and the model's defaults, can erase a country or caricature it \citep{wang2024flatten}. Both are measured and published per country. Defaulting to English prompts is a measured trade-off, not a judgement about languages; H6 tests the alternative.

\section{Conclusion}

A global agent population is many datasets of very different quality under one name. This registered report fixes, before any data are used, how the difference stays visible: provenance on every field, referent rules for copied answers, an exhaustive tier rule and a confidence score on every country, standard errors that include the donor surveys' own error, noise-corrected diagnostics for flattening and caricature, and a protocol whose primary tests can fail. The terms reading changes the design: the open synthetic skeleton rests on upstream sources that restrict commercial use, the largest attitude archive is non-profit only, and agent-level attitude vectors cannot be released. Whether the tiers and the score track real error, whether grounding moves non-Western countries away from the model's default, and whether enough post-cutoff data will exist to say so confirmatorily, are open questions. The artifacts in Table~\ref{tab:pending} will answer them.

\appendix
\section{Datasheet skeleton}
\label{app:datasheet}

Following \citet{gebru2021}, each release ships one datasheet for the resource and one row per country. \emph{Motivation:} grounding a simulated population with per-country uncertainty. \emph{Composition:} one record per agent (Section~\ref{sec:schema}); per-country $n_k$, $m_k$, frame, tier, confidence components, $q^{\text{lang}}$, validation coverage, total standard errors, excluded population share for aggregates; no agent-level attitude vectors. \emph{Collection:} derived from the sources in Table~\ref{tab:sources} at stated versions and fieldwork years; nothing collected from people by us. \emph{Preprocessing:} post-stratification or cohort shift, education crosswalk, referent classes, donor matching with recorded collapses and bootstrap draws, fallback map. \emph{Uses:} aggregate simulation with tiers shown; not for individual inference, voter targeting, or any claim of measured public opinion. \emph{Distribution:} code, schema, maps, datasheets and permitted aggregates; demographic rows only if counsel confirms; recipe for the rest. \emph{Maintenance:} versioned; any model or source change triggers the reruns in Section~\ref{sec:limits}; licence rows updated as written answers arrive; registry amendments listed.

{\raggedright
\bibliographystyle{plainnat}
\bibliography{refs}
}

\end{document}
